\documentclass[eqsecnum, authorsinfo]{kioj6}

% Пакеты, используемые непосредственно в тексте статьи.
\usepackage{amsmath}
\usepackage{booktabs}
\usepackage{tabularx}

% ==================== НАСТРОЙКИ ЖУРНАЛА ====================

\setcounter{page}{1}
% Номер выпуска и DOI не задаются авторами на этапе подачи рукописи.
\issue{2026}{--}{edu}
\journalsectionnoimage{artificial-intelligence}
\articletype{1}

\udknumber{}
\doinumber{}

% ==================== ДАННЫЕ СТАТЬИ ====================
% =============== Ru ================

\title[MPCA-Attention]
      {MPCA-Attention: многоветвистая CNN для классификации МРТ посредством ортогональных проекций}

\author{1}[Берлет М.~В.\affil{1}]{Берлет Максим Валерьевич}
\authorinfo{1}{студент 4 курса прикладной математики и информатики, СПбГЭТУ «ЛЭТИ»}
\authoremailEnvelope{1}{berletmaximq@gmail.com}
\authororcid{1}{https://orcid.org/0009-0003-0513-3806}


\author{2}[Филатов А.~Ю.\affil{1}]{Филатов Антон Юрьевич}
\authorinfo{2}{кандидат технических наук, доцент кафедры Кафедра математического обеспечения и применения ЭВМ СПбГЭТУ «ЛЭТИ»}
\authoremail{2}{aifilatov@etu.ru}
\authororcid{2}{https://orcid.org/0000-0003-4298-8523}

\affiliation{1}{Санкт-Петербургский государственный электротехнический университет «ЛЭТИ» им. В. И. Ульянова (Ленина), Российская Федерация, 197022, Санкт-Петербург,
	ул. Профессора Попова, 5}

% =============== Eng ================

\titleen[MPCA-Attention]
        {MPCA-Attention: Multi-Branch CNN for MRI Classification via Orthogonal Projections}

\authoren{1}[Berlet M.~V.\affil{1}]{Berlet Maxim}
\authorinfoen{1}{4 year student of applied math and computer science, SPBETU «LETI»}

\authoren{2}[Filatov A.~Y.\affil{1}]{Filatov Anton}
\authorinfoen{2}{PhD in Engineering, Associate Professor of the Department of Software Engineering and Computer Applications}

\affiliationen{1}{Saint Petersburg State Electrotechnical University “LETI”, Professora Popova ul. 5, St. Petersburg, 197022, Russian Federation}


\begin{document}
\russianversion
\maketitle
\pagestyle{mainstyle}

\begin{abstract}
Рассматривается задача классификации МРТ-исследований головного мозга в условиях ограниченных вычислительных ресурсов и малого объёма размеченных данных. Предложенный подход объединяет многомерный анализ главных компонент (Multilinear Principal Component Analysis, MPCA), предобученные двумерные свёрточные сети и механизм многоголового внимания. Трёхмерные МРТ-тензоры преобразуются в трёхканальные представления ортогональных анатомических плоскостей. Независимые ветви CNN извлекают признаки каждой проекции, после чего модуль позднего слияния объединяет информацию между плоскостями для классификации. Эта краткая препринт-версия представляет метод и его программную реализацию. Подробные экспериментальные протоколы и численные результаты в данной версии не приводятся.
\keywordsru{MPCA; многоветвистая свёрточная сеть; многоголовое внимание; позднее слияние; классификация МРТ.}
\end{abstract}

\section{Введение}
Объёмные МРТ-исследования содержат информацию, распределённую между различными анатомическими плоскостями. Обработка полного объёма трёхмерной сетью может требовать значительных вычислительных ресурсов, а использование отдельных срезов предполагает явный алгоритм их выбора. В работе исследуется тензорная проекция как способ получения представлений, совместимых с предобученными двумерными CNN.

Рассматриваются шесть групп: контрольная группа, болезнь Паркинсона, болезнь Альцгеймера, синдром дефицита внимания и гиперактивности, рассеянный склероз и расстройство аутистического спектра.

\section{Обзор метода}
\subsection{Предобработка и ортогональные проекции}
Исследовательский конвейер предусматривает подготовку МРТ-данных: удаление черепной коробки, выделение серого вещества, пространственную регистрацию и сглаживание. Этот внешний этап не автоматизирован Python-скриптами репозитория.

Программный конвейер загружает объёмные изображения, приводит их к общей пространственной размерности, выполняет нормализацию интенсивностей по процентильным границам и центрирование агрегированного набора данных. Алгоритм MPCA оценивает проекционные матрицы посредством итеративного обновления подпространств отдельных мод тензора. Для объёма $\mathcal{X}\in\mathbb{R}^{I_1\times I_2\times I_3}$ представление по выбранной плоскости записывается как
\begin{equation}
\mathcal{Y}^{(m)}=\mathcal{X}\times_m U_m^{\mathsf{T}},
\end{equation}
где размерность выбранной моды сокращается до трёх компонент. Аксиальная, фронтальная и сагиттальная проекции преобразуются в трёхканальные изображения. Каналы соответствуют проекционным компонентам, а не физическому цвету.

\subsection{Архитектура модели}
Три независимые ветви ConvNeXt-Tiny с предобученными весами извлекают признаки соответствующих плоскостей. Каждое изображение приводится к размеру $224\times224$ и нормализуется с использованием статистик ImageNet.

Для каждой ветви её признаки выступают запросом модуля внимания, а признаки всех трёх ветвей используются как ключи и значения. Многоголовое внимание дополняется выходной проекцией и остаточным соединением. Полученные представления конкатенируются и передаются полносвязному классификатору, формирующему логиты шести классов.

\section{Программная реализация}
Репозиторий содержит средства тензорной проекции, загрузчики изображений, многоветвистые классификаторы, взвешенную функцию потерь Focal Loss, процедуры обучения и оценки. Альтернативные реализации включают одноветвистую модель, усреднение и конкатенацию признаков, а также слияние двунаправленной рекуррентной сетью. Оценка предусматривает Accuracy, Macro F1, Macro Recall, Macro Precision и матрицу ошибок.

Настройки экспериментов задаются в Python-коде. Воспроизводимость зависит от подготовки данных, порядка файлов, индексов разбиения, версий библиотек, случайных seed, предобученных весов и устройства вычислений. При воспроизведении исследования необходимо сопоставлять экспериментальную процедуру с текущим поведением реализации.

\pagebreak

\section{Статус публикации}
Настоящий документ является сокращённой препринт-версией. Полный обзор литературы, состав выборки, экспериментальные настройки, количественные результаты и абляционное исследование здесь опущены. Документ не заменяет полную рукопись и не свидетельствует о независимом воспроизведении экспериментов.

\printauthorsinfo
\begin{translatedpart}
\maketranslatedtitle
\begin{abstract}
This work investigates brain MRI classification under limited computational resources and a small labeled dataset. The proposed approach combines Multilinear Principal Component Analysis (MPCA), pretrained two-dimensional convolutional networks, and multi-head attention. Three-dimensional MRI tensors are transformed into three-channel representations of orthogonal anatomical planes. Independent CNN branches extract features from each projection, and a late-fusion module combines information across planes for classification. This short preprint presents the method and its software implementation. Detailed experimental protocols and numerical results are omitted from this version.
\keywordsen{MPCA; multi-branch CNN; multi-head attention; late fusion; MRI classification.}
\end{abstract}
\end{translatedpart}
\end{document}
